% Introduction

Japanese animation is commercially successful and its production market is at a
record (\cref{subsec:labor}). Behind that growth is a structural problem:
the creative workforce that produces anime remains largely invisible.

An exceptionally large share of anime production labour is performed outside
employment. In the Japan Animation Creators Association's most recent survey, 37.0\% of
respondents described themselves as freelance or self-employed~\cite{janica2026survey},
about five times the share of the Japanese workforce as a whole
(\cref{subsec:labor}). They work within a multi-layered subcontracting system in which
studios let out whole episodes or individual tasks to second- and third-tier
subcontractors, and the credits that reach the audience follow the contract rather than
the work: an outsourced episode is credited to the company that produced it, and the
individuals inside it are named only sometimes.

This \textit{contribution invisibility problem} has concrete consequences. A freelance
animator cannot demonstrate a work history verifiably to a prospective client, which
weakens their position in a market where more than a third of respondents earn under
\textyen 200{,}000 a month from anime work~\cite{nafca2024survey}; and when a studio
closes, the informal record of who worked on what goes with it. Recent policy responses
reach the contract and not the credit (\cref{subsec:labor}).

In parallel, blockchain technology has introduced new primitives
for decentralized credentialing.
Weyl et al.~\cite{buterin2022decentralized} proposed \textbf{Soulbound Tokens (SBTs)}---
non-transferable tokens bound to a specific blockchain address---
as a mechanism for encoding credentials, commitments, and affiliations
without the speculative dynamics of conventional NFTs.
SBTs have since been proposed and prototyped for healthcare credentialing,
academic certificate verification, and agricultural supply chains,
but \textbf{to the best of our knowledge, no prior work has applied SBTs
to creative production workflows or to the credentialing of freelance animators.}

Existing blockchain applications in the anime industry address the audience side of the
business: IP tokenization, collectibles and fan engagement, where the revenue comes
from fans~\cite{fourpillars2025anime,animechain2024whitepaper}. This paper points the
same primitives at the labour side, where what is at stake is not revenue but a record
of who did the work.

This paper presents \textbf{AnimatorSBT}, a blockchain-based contribution
certification framework designed to integrate SBT issuance into an existing
production management tool for freelance animators.
The system is designed to create immutable, tamper-evident, and portable records
of individual contributions at the granular level of production tasks
(e.g., key animation cuts, in-between frames, coloring assignments), with
verifiability currently limited to integrity and existence
(authenticity requires the studio co-signing layer of \cref{subsec:cosigning},
which we implement as a tested reference contract but do not yet deploy
on-chain or operate with independent studio attesters).
The framework is designed to be \textit{transparent to the end user}:
animators continue their normal workflow of recording tasks in the management tool,
and contribution certificates are issued automatically
upon project milestones without requiring blockchain knowledge,
wallet management, or gas fee payment.

This paper contributes an analysis of where an attestation can come from when no
body holds jurisdiction to make one, and a measurement of what it costs to carry
that attestation on a public chain.

The analysis is that the attestation already exists. A studio that accepts a
delivery and pays for it has made a determination about who did the work, for its
own commercial reasons, recorded in a system it already keeps; a certificate
anchored to that determination records it rather than manufacturing it, which is a
different design from anchoring on a self-reported log and has different security
properties (\cref{subsec:integration,subsec:trust}). We reach it by way of a
formalisation of the contribution invisibility problem as the gap between actual
and credited contributors, and of why the institutional route that closed the same
gap in film and television is unavailable here
(\cref{subsec:formalization,subsec:gap}).

The measurements say where the cost of doing this actually falls, and it is not
where the design puts its attention: not on the chain, and not in the contract's
logic, but in the arrangement that keeps the metadata reachable
(\cref{sec:evaluation}). Both results generalise beyond this application.

Supporting these, we deploy a soulbound contract on a public testnet and provide
tested reference implementations of a studio co-signing layer and of a
deterministic evidence-hash serialisation, the latter closing a determinism gap
that made the earlier hash meaningless under key reordering
(\cref{subsec:cosigning,subsec:canonical}). We claim no novelty for the token
mechanism itself, which composes existing standards, and we identify the adoption
and authenticity problems the design does not solve (\cref{sec:discussion}).

A note on vocabulary, since this paper's argument should be followed by readers
who do not work with public ledgers. A \emph{non-transferable}, or
\emph{soulbound}, token is a record that its holder cannot sell or pass to anyone
else; to \emph{mint} one is to create it. \emph{On chain} means stored in the
public ledger itself, which every participant replicates and none can silently
alter; \emph{off chain} means stored elsewhere, with the ledger holding only a
pointer. \emph{Gas} is the unit in which a network charges for executing a
transaction, and the fee is that quantity priced in the network's token, so a cost
in gas is stable while a cost in dollars is not. A \emph{testnet} is a public
network running the same software with tokens that have no monetary value, used
for measurement rather than operation. \emph{IPFS} addresses files by a hash of
their contents rather than by location, so a file is retrievable from anyone
holding it and unretrievable when nobody does---which is what a \emph{pinning}
subscription buys. A \emph{wallet} is the key pair controlling an address, and an
address is what a certificate is issued to.
